\documentclass{pma-report}
\begin{document}
\begin{titlepage}{
    % title = {}, % необов'язковий параметр
    % type = {}, % необов'язковий параметр, "звіт про <type> роботу..." якщо title відсутній
    numero = {1}, % необов'язковий параметр
    subject = {Приколознавство},
    theme = {Використання згорткових нейронних мереж для аналізу приколів}, % необов'язковий параметр
    % variant = {}, % необов'язковий параметр
    year = {2077},
    ministry % необов'язковий параметр, вмик./вимк. перший рядок у титулці з "МОН України"
}
    Виконав:\\
    студент VII курсу групи КМ-xx,\\
    спеціальність F1\\
    Прикладна математика\\
    ПІБ\\ % за потреби місце для підпису: \signplace
    \\
    Керівник:\\
    посада, ступінь\\
    ПІБ
\end{titlepage}

\tableofcontents

\intro
Тут зазвичай наведено тему та мету роботи.

\chapter{Розділ основної частини}
Посилання на додаток --- \verb|\ref{<label>}|: "<див. додаток~\ref{app:listings}">. Тільда --- це нерозривний пробіл, щоб слово "<додаток"> і його позначення не опинилися на різних рядках.

\section{Підрозділ основної частини}
\subsection{Пункт}
\subsubsection{Підпункт (якщо дуже треба)}
\subsection{Ще один пункт}

% Приклад включення зображення vvvv
% \begin{figure}
%     \centering
%     \includegraphics{path/to/image.png}
%     \caption{Картинка із підписом.}
%     \label{fig:google-translate}
% \end{figure}

\conclusions
Тут знаходяться висновки за всю роботу.

\bibliographystyle{udstu2006}
\bibliography{bibliography}

\appendix
\chapter{Лістинги програм}\label{app:listings} % << щоб посилатись на додаток
За необхідності наводимо лістинги командою
\verb|\lstinputlisting[language=Python, caption=main.py]{main.py}|

\end{document}